\documentclass[12pt]{article}
\usepackage[T1]{fontenc}
\usepackage[utf8]{inputenc}
\usepackage{lmodern}
\usepackage{textcomp}
\usepackage{enumitem}
\usepackage{geometry}
\usepackage{graphicx}
\usepackage{amsmath, amssymb}
\geometry{margin=1in}
\begin{document}
\begin{center}
Economics - Graduate Level Assessment 13 \
Total Marks: 40 \
Answer all questions.
\end{center}

\section*{Multiple Choice (10 marks)}
\begin{enumerate}[label=\arabic*.]
\item Which of the following is a characteristic of monopolistic competition in the long run?
\begin{enumerate}[label=(\alph*)]
    \item Excess profits due to lack of competition.
    \item Firms produce identical products.
    \item High degree of interdependence among firms.
    \item Zero economic profits.
    \item Large number of buyers and one seller.
\end{enumerate}
\item How can business expectations affect the economy?
\begin{enumerate}[label=(\alph*)]
    \item Business expectations can affect the economy by affecting the level of investment.
    \item Business expectations mold the economy through fluctuations in commodity prices.
    \item Business expectations have no impact on the economy.
    \item Business expectations steer the economy by directly setting interest rates.
    \item Business expectations influence the economy through government policy changes.
\end{enumerate}
\item Which of the following is most likely to result in a shift to the right in the demand curve for orange juice?
\begin{enumerate}[label=(\alph*)]
    \item A decrease in the price of orange juice
    \item A decrease in the price of Tang
    \item An increase in the price of grapefruit juice
    \item A decrease in the income of juice drinkers
    \item A bumper crop of oranges in Florida
\end{enumerate}
\item The aggregate demand curve has a negative slope in part because when the price level increases
\begin{enumerate}[label=(\alph*)]
    \item the real quantity of money increases
    \item imports become relatively more expensive
    \item the value of cash increases
    \item exports become relatively cheaper
    \item inflation rate decreases
\end{enumerate}
\item According to Marshall and Fisher, what are thecomponents of the demand for money?
\begin{enumerate}[label=(\alph*)]
    \item Investment and Consumption
    \item Transactions Demand and Asset Demand
    \item Income and Expenditure
    \item Savings and Interest Rates
    \item Credit and Debit
\end{enumerate}
\end{enumerate}

\section*{True / False (10 marks)}
\textit{Answer True or False}
\begin{enumerate}[label=\arabic*.]
\setcounter{enumi}{5}
\item True or False: A price ceiling always results in a surplus of the good in the market due to reduced supply.
\item True or False: An increase in the price of polyester pants causes the supply of denim jeans to fall, leading to a decrease in both price and quantity.
\item True or False: Investment demand tends to increase when the stock market experiences a downturn, as investors look for safer opportunities.
\item True or False: Structural unemployment is primarily caused by fluctuations in the business cycle, rather than technological changes or industry shifts.
\item True or False: The circular flow model of pure capitalism includes all types of consumer transactions, including those between consumers.
\end{enumerate}

\section*{Long Form Response (20 marks)}
\textit{Answer each question with a clear and complete explanation.}
\begin{enumerate}[label=\arabic*.]
\setcounter{enumi}{10}
\item Discuss how anticipated inflation benefits specific institutions and individuals, analyzing the mechanisms through which they gain advantages in an inflationary environment.
\item Discuss the differences between logit and probit models in terms of their methods for transforming probabilities. How do these differences impact their application in economic analysis?
\end{enumerate}

\vfill
\noindent\textit{End of Paper}
\end{document}